\documentclass[11pt,a4paper]{article}

\usepackage[T1]{fontenc}
\usepackage[utf8]{inputenc}
\usepackage{amsmath,amssymb,amsthm}
\usepackage{mathtools}
\usepackage[margin=2.4cm]{geometry}
\usepackage{enumitem}
\usepackage{hyperref}
\usepackage{xcolor}

\hypersetup{
  colorlinks=true,
  linkcolor=blue!60!black,
  urlcolor=blue!60!black
}

\newcommand{\F}{\mathbb{F}}
\newcommand{\Z}{\mathbb{Z}}
\newcommand{\C}{\mathbb{C}}
\newcommand{\wt}{\operatorname{wt}}
\newcommand{\im}{\operatorname{Im}}
\newcommand{\ket}[1]{\lvert#1\rangle}
\newcommand{\bra}[1]{\langle#1\rvert}
\newcommand{\inner}[2]{\langle#1\mid#2\rangle}

\title{\textbf{Special topics in logic and security II\\Exam solutions --- May 16, 2024}}
\author{}
\date{}

\begin{document}
\maketitle

\section*{Problem 1 --- Weight polynomial of $C\subset\F_{11}^3$}

Let $C=\{(x,y,z)\in\F_{11}^3\mid ax+by+cz=0\}$ with $a,b,c\in\F_{11}\setminus\{0\}$ fixed.
The map $\phi:\F_{11}^3\to\F_{11}$, $\phi(x,y,z)=ax+by+cz$, is a nonzero linear functional,
so $C=\ker\phi$ is a linear subspace of dimension $n-1=2$. Thus $|C|=11^2=121$.

Let $A_i$ be the number of codewords of weight $i$.
\begin{itemize}\setlength\itemsep{0pt}
\item $A_0=1$ (only the zero word).
\item $A_1=0$: a weight-1 word, say $(x,0,0)$ with $x\neq 0$, would force $ax=0$, impossible since $a\neq 0$. The same argument applies to the other positions.
\item $A_2$: exactly one coordinate is zero. Take e.g.\ $z=0$, $x,y\neq 0$. Then $ax+by=0$ gives $y=-ab^{-1}x$. For each $x\in\F_{11}^{\times}$ this yields a unique $y\in\F_{11}^{\times}$, so $10$ codewords. By symmetry the same count holds for the cases $y=0$ and $x=0$, totalling $A_2=30$.
\item $A_3=121-1-0-30=90$.
\end{itemize}

Following the textbook convention $A_C(z)=\sum_{i=0}^{n}A_i z^i$, the weight polynomial is
\[
\boxed{\,A_C(z)\;=\;1+30\,z^{2}+90\,z^{3}\,}
\]
(equivalently, $W_C(X,Y)=X^3+30\,XY^2+90\,Y^3$ in the homogeneous form).

\section*{Problem 2 --- MacWilliams formula for $C^{\perp}$}

The MacWilliams identity for a linear $[n,k]$ code over $\F_q$ (\textit{cf.}\ codingth.pdf, \S14, Theorem~14.2) reads
\[
A^{\perp}(z)\;=\;q^{-k}\,(1+(q-1)z)^{n}\,A\!\left(\tfrac{1-z}{1+(q-1)z}\right).
\]
With $q=11$, $n=3$, $k=2$:
\[
A^{\perp}(z)=\frac{1}{121}\,(1+10z)^{3}\,A\!\left(\tfrac{1-z}{1+10z}\right).
\]
Inserting $A(z)=1+30z^{2}+90z^{3}$, the numerator $(1+10z)^{3}+30(1-z)^{2}(1+10z)+90(1-z)^{3}$ expands to
\begin{align*}
(1+10z)^{3}&=1+30z+300z^{2}+1000z^{3},\\
30(1-z)^{2}(1+10z)&=30+240z-570z^{2}+300z^{3},\\
90(1-z)^{3}&=90-270z+270z^{2}-90z^{3}.
\end{align*}
Summing the three lines: $121+0\cdot z+0\cdot z^{2}+1210\,z^{3}$. Dividing by $121$:
\[
\boxed{\,A^{\perp}(z)\;=\;1+10\,z^{3}\,}
\]

\paragraph{Interpretation.}
$C^{\perp}$ has dimension $n-k=1$, so $|C^{\perp}|=11$. It is generated by the single row $(a,b,c)$ of the check matrix of $C$. Its codewords are the scalar multiples $\lambda(a,b,c)$ for $\lambda\in\F_{11}$: one zero word (weight $0$) and ten nonzero multiples (weight $3$, since $a,b,c$ are all nonzero). This is exactly $1+10\,z^{3}$, confirming the formula.

\section*{Problem 3 --- $w=(6,6,1,2,9,0)\in G$}

The code $G$ is the classical Reed--Solomon-type evaluation code
\[
G=\bigl\{(P(1),P(2),\dots,P(6))\in\F_{11}^{6}\;\bigm|\;P\in\F_{11}[X],\;\deg P\le 2\bigr\},
\]
which has dimension $3$. To show $w\in G$ we find $P(X)=\alpha X^2+\beta X+\gamma\in\F_{11}[X]$ satisfying $w_i=P(i)$ for $i=1,\dots,6$.

Use the first three equations:
\[
\begin{cases}
\alpha+\beta+\gamma\equiv 6,\\
4\alpha+2\beta+\gamma\equiv 6,\\
9\alpha+3\beta+\gamma\equiv 1.
\end{cases}\pmod{11}
\]
Subtracting the first from the second: $3\alpha+\beta\equiv 0$.\\
Subtracting the first from the third: $8\alpha+2\beta\equiv -5\equiv 6$, i.e.\ $4\alpha+\beta\equiv 3$.\\
Subtracting these two: $\alpha\equiv 3$, hence $\beta\equiv -9\equiv 2$ and $\gamma\equiv 6-3-2=1$.

So $P(X)=3X^{2}+2X+1$. We verify on the remaining points (mod $11$):
\[
\begin{array}{c|cccccc}
i & 1 & 2 & 3 & 4 & 5 & 6\\\hline
P(i) & 6 & 17 & 34 & 57 & 86 & 121\\
P(i)\bmod 11 & 6 & 6 & 1 & 2 & 9 & 0
\end{array}
\]
All six values match $w$, so $w=(P(1),\dots,P(6))\in G$.\hfill$\blacksquare$

\section*{Problem 4 --- Which of $R_1,R_2$ is unitary?}

A real $2\times 2$ matrix $M$ is unitary iff $M^{T}M=I_2$ iff its columns are orthonormal (over the standard inner product).

\paragraph{Matrix $R_1=\tfrac{1}{2}\begin{pmatrix}1 & -\sqrt{3}\\ \sqrt{3} & 1\end{pmatrix}$.}
\[
R_1^{T}R_1=\frac{1}{4}\begin{pmatrix}1 & \sqrt{3}\\ -\sqrt{3} & 1\end{pmatrix}\begin{pmatrix}1 & -\sqrt{3}\\ \sqrt{3} & 1\end{pmatrix}=\frac{1}{4}\begin{pmatrix}1+3 & -\sqrt{3}+\sqrt{3}\\ -\sqrt{3}+\sqrt{3} & 3+1\end{pmatrix}=I_2.
\]
So $R_1$ \emph{is} unitary --- it is the rotation $R_{\pi/3}$ (since $\cos(\pi/3)=1/2$, $\sin(\pi/3)=\sqrt{3}/2$).

\paragraph{Matrix $R_2=\tfrac{1}{2}\begin{pmatrix}\sqrt{2} & -\sqrt{2}\\ \sqrt{2} & \sqrt{2}\end{pmatrix}$.}
\[
R_2^{T}R_2=\frac{1}{4}\begin{pmatrix}\sqrt{2} & \sqrt{2}\\ -\sqrt{2} & \sqrt{2}\end{pmatrix}\begin{pmatrix}\sqrt{2} & -\sqrt{2}\\ \sqrt{2} & \sqrt{2}\end{pmatrix}=\frac{1}{4}\begin{pmatrix}2+2 & -2+2\\ -2+2 & 2+2\end{pmatrix}=I_2.
\]
So $R_2$ \emph{is also} unitary --- it is the rotation $R_{\pi/4}$ (since $\tfrac{\sqrt{2}}{2}=\cos(\pi/4)=\sin(\pi/4)$).

\paragraph{Conclusion.} Both $R_1$ and $R_2$ are unitary: each is a real rotation matrix, with $\det R_i = 1$ and orthonormal columns.

\section*{Problem 5 --- $(A\otimes B)(C\otimes D)=AC\otimes BD$}

Let $A\in M_{m\times n}(K)$, $B\in M_{p\times q}(K)$, $C\in M_{n\times r}(K)$, $D\in M_{q\times s}(K)$ so that all the products below are defined.

\paragraph{Proof via simple tensors.}
The Kronecker product of two linear maps acts on a simple tensor by the rule
\[
(A\otimes B)(u\otimes v)=(Au)\otimes(Bv),
\qquad u\in K^{n},\;v\in K^{q}.
\]
Apply this twice. For any $u\in K^{r}$, $v\in K^{s}$,
\begin{align*}
(A\otimes B)(C\otimes D)(u\otimes v)
&=(A\otimes B)\bigl((Cu)\otimes(Dv)\bigr)\\
&=A(Cu)\otimes B(Dv)\\
&=(AC)u\otimes(BD)v\\
&=(AC\otimes BD)(u\otimes v).
\end{align*}
The simple tensors $\{u\otimes v\}$ span $K^{r}\otimes K^{s}$, and both sides are linear; agreement on a spanning set forces equality of the linear maps. Hence $(A\otimes B)(C\otimes D)=AC\otimes BD$.\hfill$\square$

\paragraph{Alternative: block computation.}
Writing $A\otimes B$ as the block matrix with $(i,j)$-block $a_{ij}B$ (\textit{cf.}\ quantumcomp.pdf, \S3), the $(i,k)$-block of the product equals
\[
\sum_{j=1}^{n}(a_{ij}B)(c_{jk}D)\;=\;\Bigl(\sum_{j=1}^{n}a_{ij}c_{jk}\Bigr)\,BD\;=\;(AC)_{ik}\,BD,
\]
which is exactly the $(i,k)$-block of $AC\otimes BD$.

\section*{Problem 6 --- Discrete Fourier transform on $\Z_4$}

Following the textbook convention (quantumcomp.pdf, \S7), the DFT of $f:\Z_n\to\C$ is
\[
\widehat{f}(x)=\frac{1}{\sqrt{n}}\sum_{y\in\Z_n} e^{-\frac{2\pi i\,xy}{n}}f(y).
\]
For $n=4$ we have $e^{-2\pi i/4}=e^{-i\pi/2}=-i$, so
\[
\widehat{f}(x)=\frac{1}{2}\sum_{y=0}^{3}(-i)^{xy}f(y).
\]
Writing out the four terms $y=0,1,2,3$, and using $(-i)^{2x}=(-1)^{x}$ and $(-i)^{3x}=i^{x}$:
\[
\boxed{\;\widehat{f}(x)\;=\;\tfrac{1}{2}\Bigl[\,f(0)\;+\;(-i)^{x}\,f(1)\;+\;(-1)^{x}\,f(2)\;+\;i^{x}\,f(3)\,\Bigr]\;}
\]

Evaluating at each $x\in\Z_4$:
\[
\begin{aligned}
\widehat{f}(0)&=\tfrac{1}{2}\bigl[f(0)+f(1)+f(2)+f(3)\bigr],\\
\widehat{f}(1)&=\tfrac{1}{2}\bigl[f(0)-i\,f(1)-f(2)+i\,f(3)\bigr],\\
\widehat{f}(2)&=\tfrac{1}{2}\bigl[f(0)-f(1)+f(2)-f(3)\bigr],\\
\widehat{f}(3)&=\tfrac{1}{2}\bigl[f(0)+i\,f(1)-f(2)-i\,f(3)\bigr].
\end{aligned}
\]
These four rows are exactly the rows of the unitary $4\times 4$ DFT matrix
\[
F_{4}=\frac{1}{2}\begin{pmatrix}1 & 1 & 1 & 1\\ 1 & -i & -1 & i\\ 1 & -1 & 1 & -1\\ 1 & i & -1 & -i\end{pmatrix}.
\]

\end{document}
